\chapter*{기호와 용어}
\addcontentsline{toc}{chapter}{기호와 용어}
\markboth{기호와 용어}{기호와 용어}

\renewcommand{\arraystretch}{1.2}
\begin{longtable}{>{\raggedright\arraybackslash}p{0.25\textwidth}>{\raggedright\arraybackslash}p{0.65\textwidth}}
\toprule
\textbf{기호·용어} & \textbf{뜻} \\
\midrule
\endfirsthead
\toprule
\textbf{기호·용어} & \textbf{뜻} \\
\midrule
\endhead
\bottomrule
\endfoot

$L/K$ & 체확장. \\
$[L:K]$ & $K$-벡터공간으로 본 $L$의 차원. \\
$\overline K$ & 체 $K$의 대수적 폐포. \\
분해체 & 주어진 다항식이 완전히 분해되고 모든 근으로 생성되는 최소 확장체. \\
갈루아 확장 & 정규이면서 분리인 대수적 확장. \\
$\Gal(L/K)$ & 갈루아 확장 $L/K$의 $K$-자동동형군. \\
$\Gal(f/K)$ & 분리다항식 $f$의 분해체가 갖는 $K$-갈루아군. \\
$S_n$ & $n$개 원소의 모든 순열이 이루는 대칭군. \\
$A_n$ & $S_n$의 짝순열들이 이루는 교대군. \\
$\Sym(R)$ & 집합 $R$의 모든 순열이 이루는 군. \\
$\sgn(\pi)$ & 순열 $\pi$의 부호. 짝순열이면 $1$, 홀순열이면 $-1$. \\
추이적 작용 & 임의의 두 점 사이를 어떤 군 원소가 이동시키는 군작용. \\
$G\cdot\alpha$ & 원소 $\alpha$의 $G$-궤도. \\
$\Stab_G(\alpha)$ & $\alpha$를 고정하는 $G$의 안정자 부분군. \\
판별식 $\Delta_f$ & 근 차이들의 제곱곱 $\prod_{i<j}(\alpha_i-\alpha_j)^2$. \\
우울형 삼차식 & 이차항이 없는 $X^3+aX+b$ 꼴의 삼차다항식. \\
$D_4$ & 위수 $8$인 정사각형의 대칭군. \\
$V_4$ & 클라인 네 군 $C_2\times C_2$. \\
$s_i$ & 변수들의 $i$번째 기본대칭식. \\
일반다항식 & 계수를 대수적으로 독립인 변수로 둔 $X^n+S_1X^{n-1}+\cdots+S_n$. \\
$K^{\times2}$ & $K$의 영이 아닌 제곱들의 집합. \\

\end{longtable}
